%!TEX root = ../MainBody.tex

\chapter{总结与展望}

\section{总结}

基于路径追踪的蒙特卡洛方法为耦合热输运红外仿真提供了一条无需空间网格化的求解途径，可在同一路径空间中同时处理辐射、传导、对流三类耦合热输运过程。然而，该方法在实际仿真规模下面临严峻的计算效率挑战：单条路径在完整生命周期内需执行数十万次异构几何查询，每个探测器像素需数百条独立路径以达到工程收敛精度，总查询量达数十亿量级。原始 CPU 逐路径的递归路径追踪求解架构以同步阻塞方式逐条执行几何查询，无法利用 GPU 大规模并行吞吐能力。本文针对上述问题，以 Wavefront 等价变换为核心，设计了面向 GPU 加速的并行化耦合热输运红外仿真架构，并以 CPU-GPU 混合方式完成首个工程实例化与正确性验证，包含三项主要贡献。

\paragraph{ 基于 Wavefront 架构的并行化耦合热输运求解器}（第三章）。本文提出了将递归求解器等价改造为耦合输运状态机Wavefront求解器的完整方法。首先，从并行化设计视角对原始算法进行规约分析，识别出全部 19 个几何查询挂起点、三类不可消除的查询类型（RT 射线求交、CP 最近面投影、ENC 封闭腔体成员判定）以及五类路径状态字段，为状态机设计提供完备约束。在此基础上，通过递归到状态机的等价变换，将路径状态从函数调用栈显式化为持久化 \texttt{path\_state} 结构体，解耦状态推进与几何查询执行，构建了以路径池为中心的 Wavefront 四阶段迭代执行模型（推进$\to$收集$\to$批量执行$\to$回填）。工程实现上，设计了 CPU 承担路径调度、GPU 承担批量查询的混合架构，并引入双流 Ping-Pong 异步调度机制以隐藏跨设备传输延迟。实验验证方面，物理过程一致性验证覆盖 13 个递进复杂度的验证场景，共计 282 个有效数据点全部通过 $3\sigma$ 检验；端到端统计等价性检验在 porous 场景上对 153{,}976 个有效像素执行逐像素双样本 $Z$ 检验，经 Bonferroni 和 Benjamini-Hochberg 双重多重比较校正后报告零显著差异像素，确认 Wavefront 改造不引入系统性偏差。性能实验表明，32 线程下 Wavefront求解器+OptiX 方案与递归求解器+Embree 基线达到性能均势（$1.02\times$），8\textasciitilde 16 线程典型工作站配置下实现 $1.3$\textasciitilde$1.4\times$ 端到端加速；路径池容量对比实验（8{,}192 $\to$ 16{,}384，$2.25\times$ 退化）与查询后端消融实验（WF + Embree 单池 = 426~s，Baseline 的 $4.0\times$ 退化）共同定量证明 CPU 侧 Cache Miss 为当前主导性能约束、GPU 后端是恢复性能均势的必要条件，为后续优化提供了明确方向。

\paragraph{ 基于 OptiX 的 GPU 硬件光追查询后端}（第四章）。针对 Wavefront 架构产生的三类批量查询请求，本文设计了基于 NVIDIA OptiX 的统一 GPU 查询后端。RT 查询作为硬件光追的原生操作，由 RT Core 直接加速 BVH 遍历与三角形求交，在 RTX 4090 平台上实测吞吐达 5.15 GRays/s，约为 CPU Embree 后端（i9-13900K，16 线程，15.45 MRays/s）的 333 倍。CP 查询在语义上不是射线问题——\texttt{optixTrace()} 返回射线-三角形交点而非查询点到三角形的投影点——本文通过自定义 AABB 图元代理与零长射线相结合的方法，将近接投影查询重构为等价的光追问题，实现 19.7\textasciitilde 45.9 倍 GPU/CPU 加速。ENC 查询基于 Jordan-Brouwer 分离定理，将多面体点-内外关系判定分解为一次 CP 查询与一条 $+X$ 方向射线查询的组合，仅需两次 \texttt{optixLaunch} 即完成。三类映射整合于单一 OptiX Pipeline 中，通过五套 Shader Binding Table 实现查询类型路由，后端接口与原 Embree 后端 100\% API 兼容，支持受控对比实验。

\paragraph{ 热输运场景仿真任务 GUI 编辑器}（第五章）。作为面向仿真工作流的配套工程工具，针对 Stardis 求解器纯命令行接口在场景文件手工编辑、几何缺乏可视化、表面区域手工对应、任务参数组装繁琐、多任务缺乏自动化五类系统性使用痛点，本文设计并实现了基于 Python 3 + PyQt5 + VTK 的图形化编辑器 StardisEditor。编辑器提供场景文件完整读写闭环、三维可视化与交互式画笔标注、场景树层级浏览与属性编辑、任务队列配置与自动化执行等功能，以工程文件格式统一管理场景与任务配置。Model-View 架构与信号/槽机制保证多面板间数据一致性，可执行文件标签引用机制从工作流层面联通了 Wavefront 求解器与 OptiX 后端的对比实验。

综上，本文三项贡献共同构成了从求解器内核改造、GPU 查询后端设计到用户工具链建设的完整技术体系。Wavefront 架构实现了路径状态与几何查询的解耦，为批量化 GPU 加速提供了架构基础；OptiX 后端将三类异构查询统一映射至 GPU 硬件光追管线，在关键查询类型上实现了两到三个数量级的吞吐提升；GUI 编辑器作为配套工程工具将上述算法能力封装为可重复、低门槛的工程实践。物理一致性验证与端到端统计等价性检验共同确认，混合架构在实现性能提升的同时不引入系统性物理偏差。上述成果有效验证了 Wavefront 架构的正确性与查询后端的吞吐能力，同时精确定位了 CPU 侧实现的性能天花板，为全 GPU 架构演进提供了明确的实验依据。系统化的几何敏感性与超大规模场景扩展性验证留待后续工作。

\section{工作展望}

尽管本文在 Wavefront 并行化架构、GPU 光追后端和 GUI 工具链三个方面取得了阶段性成果，当前系统仍存在若干局限性，可从以下方向展开后续研究。

\subsection{ 全 GPU 求解器架构 }当前混合架构中，路径状态调度（advance/cascade/distribute）保留在 CPU 侧执行，$T_A/T_B$ 时间分解（表~3.6）显示 GPU 设备侧仅占墙钟时间的 10.9\%，端到端加速上界仅为 $1.12\times$；路径池容量从 8{,}192 增至 16{,}384 即引发 $2.25\times$ 性能退化（第三章 3.4.3）。消除该瓶颈的根本途径是将路径状态调度整体迁至 GPU ，并在 GPU 端原生组织 Wavefront 执行模型。全 GPU 求解还可消除显式 \texttt{optixLaunch} 在每轮迭代间产生的 host 端 kernel 启动开销；查询经由 OptiX 内联接口（\texttt{optixTraverse} / \texttt{optixReportIntersection}）在调度 kernel 内直接发起，使 FSM 状态推进与查询执行融合为同一指令流。

全 GPU 路线的核心工程难点在于 GPU 端调度内核的实现。调度内核主要由 Wavefront 、FSM 和更底层的场景物理信息访问函数组成。在 Wavefront 层，调度本身不再是 CPU 主控的批次循环，而是 GPU 上的持久线程网格：每个 warp 持续从路径池消费 active 路径，完成“状态推进 $\to$ 查询发起 $\to$ 结果写回”的全过程；由于 FSM 的数十个状态在同一 warp 的 32 个 lane 内极可能高度分歧、直接执行会显著拉低 SIMT 利用率，可借鉴现代高性能 Wavefront 路径追踪器的成熟做法，在每轮推进前以 stream compaction 按状态 ID 对 active 路径重排，使 warp 内 lane 状态同质化以提高一致性。在 FSM 层，本文混合架构所定义的 \texttt{path\_state} 需进一步推进 SoA 化，使 advance 阶段对热字段的访问由 lane 级分散读写转换为 warp 级合并访存；由于 GPU 实现无需在几何查询处挂起，还可将挂起点前后的 FSM 状态融合并将几何查询直接内联，显著减少状态数与跨状态数据搬运，进一步降低线程分歧；此外，本文方法对 PicardN 迭代的求解采用了深度受限显式栈，这在 GPU 端并不适用，但可通过将具有树状关系的路径状态展平为线性结构并以 parent index 进行关联的方式实现等价功能，相关技术细节留待后续工作展开。针对更底层的场景物理信息访问函数，GPU 端实现的主要挑战在于构造与基于函数指针的 CPU 实现等价的高性能抽象接口，以支持多样化的场景几何与热物理模型；经过对这类函数的分析，其遵循统一的调用规约、访问模式高度数据驱动，可直接从绑定在几何上的物理属性数据中读取所需信息。

\subsection{ 借鉴现代高性能路径追踪的工程优化技术 }本文 Wavefront 求解器以耦合热输运的物理正确性为首要目标完成实例化，在调度组织与算法层面仍保有可观的提升空间。现代图形界的高性能路径追踪器（如 NVIDIA Falcor / NRD 体系、商用 DXR 渲染管线）所积累的大量工程优化与本文求解器结构高度同源，可按两条主线在前述全 GPU 实例化基础上有选择地引入。

第一条主线是\emph{warp 同质化}，本质是将“行为相近的路径聚到同一 warp”做到极致。除前述按状态 ID 的 stream compaction 之外，可按事件类型将每轮路径推进进一步拆分为多个独立 kernel（生成、扩展、阴影、交互、终止），使每个 kernel 内 lane 同质性更高；在材质评估阶段对路径按材质 ID 或查询类型做二次排序（Coherent Shading）使同 warp 内 lane 访问相同材质资源；并借助 NVIDIA Shader Execution Reordering（SER）由硬件在 OptiX 内联调用点自动完成 warp 重组。该主线与本文 FSM 的多状态结构天然契合，是前述 Wavefront 层调度内核的延伸优化。

第二条主线是\emph{路径剪枝}，本质是“按贡献裁减低收益路径”。蒙特卡洛路径的能量贡献存在显著长尾，Russian Roulette 与自适应路径长度可在期望意义下保持无偏的同时大幅截断低贡献路径段。该机制与本文耦合热输运框架天然兼容，与调度架构正交，可单独引入并预期在 PicardN 高阶展开与多次散射场景下带来直接收益。

上述两条主线与本文 FSM、\texttt{path\_state}、批量查询接口设计正交，无需对架构作根本性改动即可逐项引入。

\subsection{ 更丰富的物理模型 }本文采用漫射灰体、透明非参与流体与线性化对流系数三条工程简化（见第二章末），覆盖了主流工业红外仿真任务。当假设失效时，所需扩展同时涉及学术与工程两个层面：学术层面关注物理模型在路径空间中的可表达性、蒙特卡洛估计器的无偏性与方差可控性；工程层面关注 FSM 与查询管线的结构改动。两个层面缺一不可。

\emph{参与性介质失效}（烟气、火焰、薄云、气溶胶）。学术上，路径空间从"表面—表面"扩展到"表面—体—表面"，需在路径密度中引入沿光程的指数衰减项 $\exp(-\!\int_0^s \sigma_t(\ell)\,\mathrm d\ell)$ 与散射相函数 $f_p$。估计器无偏性要求自由飞行距离采样严格匹配介质消光系数（如 Woodcock null-collision 跟踪）；散射事件需通过相函数的解析或表式重要性采样保持概率密度归一；多次散射路径的方差控制依赖 Russian Roulette 与 Next-Event Estimation 等成熟方差缩减技术。Bati 等\cite{batiCouplingConductionConvection2023} 单路径框架在原理上已涵盖该方向，但 PicardN 在非线性辐射边界与体积参与项耦合下的收敛性与无偏性需重新分析证明。工程上，FSM 需新增"介质内自由飞行—散射/吸收事件"段落（约 5\textasciitilde 7 个状态），GPU 查询管线需扩展 RT 类型以支持体积光线步进与采样；该类扩展属于结构性改动，但与本文 RT/CP/ENC 三类查询保持正交。

\emph{多光谱辐射失效}（强温度依赖发射率、谱选择性材料、大气吸收带）。学术上，单波段灰体假设失效后，能量沿路径不再由单一频道携带，需在路径空间中引入波长维度；常见处理方式有按光谱发射密度采样的 Spectral Importance Sampling 与 Hero Wavelength + 重要性采样方案，两者在统计上均可保持无偏但方差—成本权衡不同；$\varepsilon(\lambda,T)$ 与 Planck 函数 $B_\lambda(T)$ 在 PicardN 内展开的一致性条件需重新推导。工程上，在 \texttt{path\_state} 中扩展波段索引字段，材质系统由灰体单值扩展为按谱段查表，采样阶段引入波段抽样；FSM 状态拓扑与查询管线本身无需结构性改动。

\emph{强非线性对流系数失效}（自然对流主导、相变换热等）。学术上，线性化对流系数 $h_{cv}=\mathrm{const}$ 在 PicardN 的 null-collision 分支选择中被显式使用以保证概率归一与无偏；若 $h_{cv}=h_{cv}(T)$ 强非线性，需将其作为 PicardN 框架内的非线性项与辐射项 $h_{rd}(T)$ 同等处理，重新推导多温度采样的乘积估计形式；当 $h_{cv}$ 跨越数量级时，估计器方差的可控性需通过 majorant 上界或额外方差缩减手段保证。工程上，在 advance 阶段引入温度依赖的换热系数闭包即可，FSM 与查询管线均保持不变。

总体而言，三类假设失效在学术上的挑战远大于工程上的改动。在学术层面问题得到解决的前提下，工程层面的扩展可在本文 \texttt{path\_state} + 批量查询接口范式内逐步实施。上述扩展的理论分析、原型实现与无偏性验证留待后续工作。

\subsection{ 仿真管线与工业/DCC 软件的数据互通 }当仿真目标从单一构件级算例扩展到整机级或系统级工业场景时，几何资产规模与边界条件复杂度迅速膨胀，单一 GUI 编辑器以手工方式组织几何、材质、边界条件、传感器布置与任务参数已不可持续——编辑管线本身成为新的瓶颈。务实的扩展方向是将本文求解器与编辑器纳入现有工业 / DCC 工具链的数据流，使 StardisEditor 在两个职能维度上从“全功能编辑器”退守为“红外仿真求解 + 结果可视化”节点。

\emph{静态场景数据维度}：将本文 BNF 文法的场景描述以更通用的结构化格式（JSON Schema、USD Schema 等）重表达，作为统一的入口数据契约；其上层资产数据通过中性格式与工业/DCC 生态对接——CAD/CAE 侧由 STEP/IGES、CGNS 等中性格式承担装配级几何与已划分网格的导入，DCC 侧由 USD、glTF、Alembic 等开放资产格式与 Blender、Maya、Houdini 等工具共享几何与光学/热学材质属性。这一维度的核心是让 StardisEditor 退出“建模与材质指派”的职能。

\emph{动态任务流程维度}：把参数扫掠、批处理、不确定度分析等大规模算例的编排能力让渡给更成熟的工作流引擎（如 ParaView、Cassiopée），StardisEditor 仅作为求解任务节点参与编排。这一维度的核心是让 StardisEditor 退出“任务调度”的职能。

在上述两维互通的基础上，交互式预览可作为面向终端用户体验的附属增强保留：结合基于深度学习的蒙特卡洛降噪算法，以低 SPP（如 1\textasciitilde 4）配合实时降噪实现秒级红外预览，使用户在跨工具编辑场景参数时即时获得温度场反馈。